% GENERATED FILE. Edit canonical lessons, then run npm run generate:books.
% canonical-source-hash: fnv1a64:bb6069fdd0812b03
% canonical-lessons: KA-C493-sittagu, KA-C493-kenaku, KA-C493-hiyalisu, KA-C493-hatase, KA-C493-udvega

\chapter{Anger and Mockery}
\label{ch:ka-anger-and-mockery}
\hlchaptermodality{493}

\begin{chapteropening}
\textbf{By the end of this chapter:} \emph{I can say to get annoyed, to provoke, to ridicule, frustration and agitation.}

Complete the last lesson of chapter 493: I can say to get annoyed, to provoke, to ridicule, frustration and agitation.
\end{chapteropening}

\section[\texorpdfstring{\kn{ಸಿಟ್ಟಾಗು} (siṭṭāgu)}{ಸಿಟ್ಟಾಗು (siṭṭāgu)}]{\kn{ಸಿಟ್ಟಾಗು} (siṭṭāgu) — to get annoyed; to lose your temper}

\label{lesson:KA-C493-sittagu}



\begin{quote}
Before the new one: say the Kannada for delight, then the Kannada for satisfaction.
\end{quote}

\subsection*{You'll want to know: \kn{ಸಿಟ್ಟಾಗು}}
\textbf{\kn{ಸಿಟ್ಟಾಗು}} — \emph{siṭṭāgu} — \textquotedblleft{}to get annoyed; to lose your temper\textquotedblright{}.

\begin{quote}
\kn{ತಡವಾಗಿ} \kn{ಬಂದಿದ್ದಕ್ಕೆ} \kn{ಅಪ್ಪನಿಗೆ} \kn{ನನ್ನ} \kn{ಮೇಲೆ} \kn{ಸಿಟ್ಟಾಯಿತು}.
\end{quote}

\begin{quote}
\emph{taḍavāgi bandiddakke appanige nanna mēle siṭṭāyitu.}
\end{quote}

\begin{quote}
Father got annoyed with me for coming late.
\end{quote}

Past (it …): \kn{ಸಿಟ್ಟಾಯಿತು}, \emph{siṭṭāyitu}.

The one who gets annoyed takes -\kn{ಗೆ}, the target … \kn{ಮೇಲೆ}. \kn{ಸಿಟ್ಟು}, a temper, is the everyday word beside \kn{ಕೋಪ}.

\begin{cousinweb}
\textbf{\kn{ಸಿಟ್ಟಾಗು}} (\emph{siṭṭāgu}) is built from \emph{siṭṭu}, \textquotedblleft{}a temper, anger\textquotedblright{}, and \kn{ಆಗು}, \textquotedblleft{}to become\textquotedblright{}; the origin of \emph{siṭṭu} is not settled.
\end{cousinweb}

\begin{grammarlens}[title={What you've built}]
One more word for describing what you lived and felt.
\end{grammarlens}

\subsection*{Your turn}
\begin{itemize}
\raggedright
  \item \emph{Say it:} \emph{siṭṭāgu}
  \item \emph{Say it:} \emph{siṭṭāgu}, once more
  \item \emph{Say it:} \emph{taḍavāgi bandiddakke appanige nanna mēle siṭṭāyitu.} — Father got annoyed with me for coming late.
  \item \emph{Recall:} say the Kannada for delight, then the Kannada for satisfaction, then say \emph{siṭṭāgu} again
\end{itemize}

\begin{tcolorbox}[breakable,colback=teal!4,colframe=teal!35!black,title={Before you move on}]
What does \kn{ಸಿಟ್ಟಾಗು} mean? (\textquotedblleft{}To get annoyed; to lose your temper\textquotedblright{}.) Say it once more.
\end{tcolorbox}
\section[\texorpdfstring{\kn{ಕೆಣಕು} (keṇaku)}{ಕೆಣಕು (keṇaku)}]{\kn{ಕೆಣಕು} (keṇaku) — to provoke; to needle}

\label{lesson:KA-C493-kenaku}



\begin{quote}
Before the new one: say the Kannada for satisfaction, then the Kannada for to get annoyed.
\end{quote}

\subsection*{You'll want to know: \kn{ಕೆಣಕು}}
\textbf{\kn{ಕೆಣಕು}} — \emph{keṇaku} — \textquotedblleft{}to provoke; to needle\textquotedblright{}.

\begin{quote}
\kn{ಅಣ್ಣ} \kn{ತಂಗಿಯನ್ನು} \kn{ಕೆಣಕಿದಾಗ} \kn{ಅವಳು} \kn{ಅತ್ತಳು}.
\end{quote}

\begin{quote}
\emph{aṇṇa taṅgiyannu keṇakidāga avaḷu attaḷu.}
\end{quote}

\begin{quote}
When the elder brother needled his younger sister, she cried.
\end{quote}

Past (I …): \kn{ಕೆಣಕಿದೆ}, \emph{keṇakide}.

The one provoked takes -\kn{ಅನ್ನು}. Close to \kn{ರೇಗಿಸು}, but meaner.

\begin{cousinweb}
\textbf{\kn{ಕೆಣಕು}} (\emph{keṇaku}) is native Dravidian: DEDR 1944 lists Kannada \emph{keṇaku}, \textquotedblleft{}to irritate, provoke, by words or deeds\textquotedblright{}.
\end{cousinweb}

\begin{grammarlens}[title={What you've built}]
One more word for describing what you lived and felt.
\end{grammarlens}

\subsection*{Your turn}
\begin{itemize}
\raggedright
  \item \emph{Say it:} \emph{keṇaku}
  \item \emph{Say it:} \emph{keṇaku}, once more
  \item \emph{Say it:} \emph{aṇṇa taṅgiyannu keṇakidāga avaḷu attaḷu.} — When the elder brother needled his younger sister, she cried.
  \item \emph{Recall:} say the Kannada for satisfaction, then the Kannada for to get annoyed, then say \emph{keṇaku} again
\end{itemize}

\begin{tcolorbox}[breakable,colback=teal!4,colframe=teal!35!black,title={Before you move on}]
What does \kn{ಕೆಣಕು} mean? (\textquotedblleft{}To provoke; to needle\textquotedblright{}.) Say it once more.
\end{tcolorbox}
\section[\texorpdfstring{\kn{ಹೀಯಾಳಿಸು} (hīyāḷisu)}{ಹೀಯಾಳಿಸು (hīyāḷisu)}]{\kn{ಹೀಯಾಳಿಸು} (hīyāḷisu) — to ridicule; to taunt}

\label{lesson:KA-C493-hiyalisu}



\begin{quote}
Before the new one: say the Kannada for to get annoyed, then the Kannada for to provoke.
\end{quote}

\subsection*{You'll want to know: \kn{ಹೀಯಾಳಿಸು}}
\textbf{\kn{ಹೀಯಾಳಿಸು}} — \emph{hīyāḷisu} — \textquotedblleft{}to ridicule; to taunt\textquotedblright{}.

\begin{quote}
\kn{ಅವನ} \kn{ಹಳೆಯ} \kn{ಬಟ್ಟೆ} \kn{ನೋಡಿ} \kn{ಕೆಲವರು} \kn{ಹೀಯಾಳಿಸಿದರು}.
\end{quote}

\begin{quote}
\emph{avana haḷeya baṭṭe nōḍi kelavaru hīyāḷisidaru.}
\end{quote}

\begin{quote}
Some people ridiculed him for his old clothes.
\end{quote}

Past (I …): \kn{ಹೀಯಾಳಿಸಿದೆ}, \emph{hīyāḷiside}.

The one ridiculed takes -\kn{ಅನ್ನು}. \kn{ಹೀಯಾಳಿಕೆ}, ridicule.

\begin{cousinweb}
\textbf{\kn{ಹೀಯಾಳಿಸು}} (\emph{hīyāḷisu}) is a native Kannada verb, \textquotedblleft{}to ridicule, to taunt\textquotedblright{}; its origin is not settled.
\end{cousinweb}

\begin{grammarlens}[title={What you've built}]
One more word for describing what you lived and felt.
\end{grammarlens}

\subsection*{Your turn}
\begin{itemize}
\raggedright
  \item \emph{Say it:} \emph{hīyāḷisu}
  \item \emph{Say it:} \emph{hīyāḷisu}, once more
  \item \emph{Say it:} \emph{avana haḷeya baṭṭe nōḍi kelavaru hīyāḷisidaru.} — Some people ridiculed him for his old clothes.
  \item \emph{Recall:} say the Kannada for to get annoyed, then the Kannada for to provoke, then say \emph{hīyāḷisu} again
\end{itemize}

\begin{tcolorbox}[breakable,colback=teal!4,colframe=teal!35!black,title={Before you move on}]
What does \kn{ಹೀಯಾಳಿಸು} mean? (\textquotedblleft{}To ridicule; to taunt\textquotedblright{}.) Say it once more.
\end{tcolorbox}
\section[\texorpdfstring{\kn{ಹತಾಶೆ} (hatāśe)}{ಹತಾಶೆ (hatāśe)}]{\kn{ಹತಾಶೆ} (hatāśe) — frustration; despair}

\label{lesson:KA-C493-hatase}



\begin{quote}
Before the new one: say the Kannada for to provoke, then the Kannada for to ridicule.
\end{quote}

\subsection*{You'll want to know: \kn{ಹತಾಶೆ}}
\textbf{\kn{ಹತಾಶೆ}} — \emph{hatāśe} — \textquotedblleft{}frustration; despair\textquotedblright{}.

\begin{quote}
\kn{ಮೂರು} \kn{ಸಲ} \kn{ಸೋತ} \kn{ಮೇಲೆ} \kn{ಅವನಿಗೆ} \kn{ಹತಾಶೆ} \kn{ಆಯಿತು}.
\end{quote}

\begin{quote}
\emph{mūru sala sōta mēle avanige hatāśe āyitu.}
\end{quote}

\begin{quote}
After losing three times, he felt frustrated.
\end{quote}

From \kn{ಹತ}, struck down, and \kn{ಆಶೆ}, hope: a hope that has been knocked down.

\begin{cousinweb}
\textbf{\kn{ಹತಾಶೆ}} (\emph{hatāśe}) is built from Sanskrit \emph{hatāśa}, \textquotedblleft{}whose hopes are destroyed, desperate\textquotedblright{} (Monier-Williams): \emph{hata}, \textquotedblleft{}struck\textquotedblright{}, and \emph{āśā}, \textquotedblleft{}hope\textquotedblright{}.
\end{cousinweb}

\begin{grammarlens}[title={What you've built}]
One more word for describing what you lived and felt.
\end{grammarlens}

\subsection*{Your turn}
\begin{itemize}
\raggedright
  \item \emph{Say it:} \emph{hatāśe}
  \item \emph{Say it:} \emph{hatāśe}, once more
  \item \emph{Say it:} \emph{mūru sala sōta mēle avanige hatāśe āyitu.} — After losing three times, he felt frustrated.
  \item \emph{Recall:} say the Kannada for to provoke, then the Kannada for to ridicule, then say \emph{hatāśe} again
\end{itemize}

\begin{tcolorbox}[breakable,colback=teal!4,colframe=teal!35!black,title={Before you move on}]
What does \kn{ಹತಾಶೆ} mean? (\textquotedblleft{}Frustration; despair\textquotedblright{}.) Say it once more.
\end{tcolorbox}
\section[\texorpdfstring{\kn{ಉದ್ವೇಗ} (udvēga)}{ಉದ್ವೇಗ (udvēga)}]{\kn{ಉದ್ವೇಗ} (udvēga) — agitation; excitement}

\label{lesson:KA-C493-udvega}



\begin{quote}
Before the new one: say the Kannada for to ridicule, then the Kannada for frustration.
\end{quote}

\subsection*{You'll want to know: \kn{ಉದ್ವೇಗ}}
\textbf{\kn{ಉದ್ವೇಗ}} — \emph{udvēga} — \textquotedblleft{}agitation; excitement\textquotedblright{}.

\begin{quote}
\kn{ಉದ್ವೇಗದಲ್ಲಿ} \kn{ಅವನು} \kn{ಜೋರಾಗಿ} \kn{ಮಾತಾಡಿದನು}.
\end{quote}

\begin{quote}
\emph{udvēgadalli avanu jōrāgi mātāḍidanu.}
\end{quote}

\begin{quote}
In his agitation he spoke loudly.
\end{quote}

\kn{ಉದ್ವೇಗಕ್ಕೆ} \kn{ಒಳಗಾಗು}, to get worked up.

\begin{cousinweb}
\textbf{\kn{ಉದ್ವೇಗ}} (\emph{udvēga}) is borrowed from Sanskrit \emph{udvega}, \textquotedblleft{}trembling, shaking; agitation, anxiety\textquotedblright{} (Monier-Williams), from \emph{ud-}, \textquotedblleft{}up\textquotedblright{}, and \emph{vega}, \textquotedblleft{}force\textquotedblright{}.
\end{cousinweb}

\begin{grammarlens}[title={What you've built}]
One more word for describing what you lived and felt.
\end{grammarlens}

\subsection*{Your turn}
\begin{itemize}
\raggedright
  \item \emph{Say it:} \emph{udvēga}
  \item \emph{Say it:} \emph{udvēga}, once more
  \item \emph{Say it:} \emph{udvēgadalli avanu jōrāgi mātāḍidanu.} — In his agitation he spoke loudly.
  \item \emph{Recall:} say the Kannada for to ridicule, then the Kannada for frustration, then say \emph{udvēga} again
\end{itemize}

\begin{tcolorbox}[breakable,colback=teal!4,colframe=teal!35!black,title={Before you move on}]
What does \kn{ಉದ್ವೇಗ} mean? (\textquotedblleft{}Agitation; excitement\textquotedblright{}.) Say it once more.
\end{tcolorbox}
